\documentclass{article}
\newcommand{\dom}{\mathrm{dom}}
\usepackage{lineno}
\linenumbers

% if you need to pass options to natbib, use, e.g.:
%     \PassOptionsToPackage{numbers, compress}{natbib}
% before loading neurips_2020

% ready for submission
% \usepackage{neurips_2020}

% to compile a preprint version, e.g., for submission to arXiv, add add the
% [preprint] option:
%     \usepackage[preprint]{neurips_2020}

% to compile a camera-ready version, add the [final] option, e.g.:
%     \usepackage[final]{neurips_2020}

% to avoid loading the natbib package, add option nonatbib:
%     \usepackage[nonatbib]{neurips_2020}

\usepackage[utf8]{inputenc} % allow utf-8 input
\usepackage[T1]{fontenc}    % use 8-bit T1 fonts
\usepackage{hyperref}       % hyperlinks
\usepackage{url}            % simple URL typesetting
\usepackage{booktabs}       % professional-quality tables
\usepackage{amsfonts}       % blackboard math symbols
\usepackage{nicefrac}       % compact symbols for 1/2, etc.
\usepackage{microtype}      % microtypography
\usepackage{amssymb,amsmath,amsthm}
\usepackage[framemethod=tikz]{mdframed}
\usepackage{verbatim}
\usepackage{bbm}

%\usepackage{mystyle}
\usepackage{mathtools,booktabs,tabu,graphicx,hyperref,epstopdf,url,rotating}
%\usepackage{fullpage}
\usepackage{algorithm}
\usepackage{cite}
\usepackage{algorithmic}
\usepackage{fancyref}
\usepackage[framemethod=tikz]{mdframed}

\newtheorem{theorem}{Theorem}[section]
\newtheorem{corollary}{Corollary}[theorem]	
\newtheorem{condition}{Condition}[theorem]	
\newtheorem{proposition}{Proposition}[theorem]
\newtheorem{definition}[theorem]{Definition}
\newtheorem{example}[theorem]{Example}
\newtheorem{assumption}[theorem]{Assumption}
\newtheorem{lemma}[theorem]{Lemma}
\newtheorem{remark}{Remark}[theorem]

\newcommand{\tstep}{{\Delta t}}

\renewcommand{\vec}[1]{{\mathbf{#1}}} 
\newcommand{\Oo}{\mathcal{O}}
\newcommand{\Nn}{\mathcal{N}}
\newcommand{\Id}{\mathrm{Id}}
\newcommand{\todolist}[1]{\textcolor{red}{(TODO: #1)}}
\newcommand{\eqdef}{\coloneq}
\newcommand{\Prox}{\mathrm{Prox}}

\newcommand{\pa}[1]{\pp{#1}}
\newcommand{\PP}{\mathbb{P}}
\newcommand{\EE}{\mathbb{E}}
\newcommand{\RR}{\mathbb{R}}

\newcommand{\clip}{L}
\newcommand{\clower}{K}
\newcommand{\cinfty}{B}
\newcommand{\ccc}{\clower}

%\usepackage{graphicx}
%\usepackage{subfigure}
\usepackage{subcaption}
\usepackage{geometry}
\geometry{
 a4paper,
 total={170mm,257mm},
 left=20mm,
 top=20mm,
 }

\DeclareMathOperator*{\diag}{diag}
\DeclareMathOperator*{\argmax}{arg\,max\,}
\DeclareMathOperator*{\argmin}{arg\,min\,}
\DeclareMathOperator*{\tr}{tr}


\title{Cartesian Parametrisation of Hadamard Langevin and Splitting Method}
% The \author macro works with any number of authors. There are two commands
% used to separate the names and addresses of multiple authors: \And and \AND.
%
% Using \And between authors leaves it to LaTeX to determine where to break the
% lines. Using \AND forces a line break at that point. So, if LaTeX puts 3 of 4
% authors names on the first line, and the last on the second line, try using
% \AND instead of \And before the third author name.

\author{%
 Ivan Cheltsov
}

\newcommand{\OO}{\mathbb{O}}

\newcommand{\commentout}[1]{}


\begin{document}

\maketitle
\section{Introduction}
In a deterministic framework, inverse problems are typically formulated as optimization problems where the goal is to find the best estimate of an unknown parameter or function based on observed data. Given a forward model $F$ that relates the unknown $x$ to the observed data $y$ through 
\begin{equation} \label{lasso_MAP}
y=F(x) + \eta \\
\end{equation}
the inverse problem seeks to recover $x$ from $y$. Since inverse problems are often ill-posed, meaning solutions may not exist, be non-unique, or be sensitive to perturbations in data, a common approach is to define a cost functional $C(x)$ that quantifies the discrepancy between the observed data and the model prediction, often incorporating regularization to handle ill-posedness. One regularization approach of interest to us is Lasso regularization, which incorporates prior knowledge about the sparsity of solution $x$ through an extra penalty term. For instance, a solution for a Lasso-regularized least squares model looks like
\begin{equation} \label{lasso_MAP}
\hat{x} = \argmin_{x} C(x) +\lambda R(x), \\
\end{equation}
with cost $C(x) = \frac{1}{2}\|F(x) - y\|^2$ and a regularization $R(x) = \|x\|_1$, weighted by a parameter $\lambda$. Other regularization techniques, such as total variation (TV) regularization, are commonly used in applications like image reconstruction and compressed sensing. Deterministic methods provide a single "best" solution but do not quantify uncertainty in the reconstruction.

Bayesian methods interpret inverse problems probabilistically by treating the unknown $x$ as a random variable with a prior distribution that encodes prior beliefs about its properties. Observed data is incorporated through a likelihood function, and Bayes’ theorem is used to compute the posterior distribution:
\begin{equation} \label{lasso_bayes}
p(x | y) \propto p(y | x)p(x), \\
\end{equation}
where $p(x | y)$ is the likelihood, $p(y | x) \propto \exp(-C(x))$ and $p(x) \propto \exp(-\lambda R(x))$ is the prior. The posterior distribution $p(x | y)$ represents the updated knowledge about $x$ after observing $y$. The most straightforward way to use the posterior is to compute the maximum a posteriori estimate (MAP) $x_{MAP} = \argmax_{x} p(x | y)$, namely, the solution of (\ref{lasso_MAP}).

A key advantage of the Bayesian approach is that it naturally quantifies uncertainty in the solution. Instead of producing a single estimate, it provides a distribution over possible solutions, allowing for credible intervals and uncertainty assessments.

\section{Sparse Priors}

\section{Setup}
We can use Langevin sampling to produce samples from $\pi$. Write $\pi\propto\exp(-\beta\, H)$ for $H((u,v))=\frac 12 \,\lambda\, \norm{(u,v)}^2 + G(u\odot v)- \frac{1}{\beta}\,\sum_i\log u_i$. The gradient
\[
\nabla H((u,v))=\lambda\begin{pmatrix}
    u\\v
\end{pmatrix}
+\begin{pmatrix}
    v\\ u
\end{pmatrix}\odot \nabla G(u\odot v)
-\begin{pmatrix}
    1/\beta u\\ 0
\end{pmatrix},
\]
where we interpret $1/u\in\real^d$ as the vector with entries $1/u_i$. Therefore, the corresponding Langevin equations are 
\begin{equation} \label{uvproc}
\begin{split}
du_t&=-\lambda\,   u_t\,dt - v_t\odot \nabla G(u_t\odot v_t)\,dt
 +\frac{1}{\beta \,u_t}\,dt+ \sqrt\frac{2}{\beta}\,dW^1_t,\\
dv_t&=-\lambda\,  v_t\,dt -  \,u_t \odot \nabla G(u_t\odot v_t)\,dt+  \sqrt\frac{2}{\beta}\,dW^2_t,
\label{uvproc}
\end{split}
\end{equation}

We can reparametrize into cartesian coordinates to obtain the system

\begin{equation} 
\begin{split} 
dy_t&=-\lambda\,   y_t\,dt - \frac{y_t}{u_t} \odot v_t\odot \nabla G(x_t)\,dt
 +\sqrt\frac{2}{\beta}\,dW^0_t,\\
dz_t&=-\lambda\,   z_t\,dt - \frac{z_t}{u_t} \odot v_t\odot \nabla G(x_t)\,dt
 +\sqrt\frac{2}{\beta}\,dW^1_t,\\
dv_t&=-\lambda\,   v_t\,dt - u_t\odot \nabla G(x_t)\,dt
 +\sqrt\frac{2}{\beta}\,dW^2_t
\label{yzvproc}
\end{split}
\end{equation}
with $y,z,v \in \mathbb{R}^d$, $\mathbb{R}^{d}$-valued Brownian motion $W^0(t), W^1(t), W^2(t)$ and
\[u_t = (y_t \odot y_t + z_t \odot z_t)^{1/2}, \quad x_t = u_t \odot v_t.\]
More concisely, we can rewrite the above as 
\begin{equation}  
 \begin{aligned}
 d\vec x= \Bigl[-\lambda \vec x - \nabla G_{3}(\vec x)\Bigr]\,dt +\sqrt{2/\beta}\, dW^3(t),
  \end{aligned}\label{cyl}
\end{equation}
for $\vec x \in \mathbb{R}^{3d}$, $G_{3}(\vec x)= G( (y\odot y + z\odot z)^{1/2} \odot v)$ and an $\mathbb{R}^{3d}$-valued Brownian motion $W^3(t)$. 

We can show, via It\^{o}'s Formula with $u_t = (y_t \odot y_t + z_t \odot z_t)^{1/2}$, that 

$$
\begin{aligned}
  du_t&=\frac{1}{u_t}\odot(y_t\odot(-\lambda\,y_t)+z_t\odot(-\lambda\, z_t))\,dt\\&\quad
  -\frac{1}{u_t}(y_t\odot (y_t/u_t)+z_t\odot(z_t/u_t))\odot v_t\odot \nabla G(u_t\odot v_t)\,dt \\&\quad+\sqrt\frac{2}{\beta} \frac 1{u_t} \odot\begin{pmatrix}y_t\\z_t\end{pmatrix}
  \cdot\begin{pmatrix} dW_t^{3,y}\\ dW_t^{3,z}\end{pmatrix}
  + \frac 12 \frac{2}{\beta} \left(\frac{-1}{u^3}\odot(y_t\odot y_t+z_t\odot z_t)+\frac{2}{u_t}\right)\,dt\\
  &=-\lambda\,   u_t\,dt - v_t\odot \nabla G(u_t\odot v_t)\,dt
 +\frac{1}{\beta \,u_t}\,dt+ \sqrt\frac{2}{\beta}\,dW_t,
\end{aligned}
$$
hence the $(u,v)$-process (\ref{uvproc}) arises from the solution of the $(y,z,v)$-process (\ref{yzvproc}). By Proposition (include reference to paper or Proposition with proof), the process (\ref{uvproc}) also inherits properties from (\ref{yzvproc}) such as being well-posed, strong Feller, irreducible. The processes agree as long as the initial distributions agree, i.e., that the probability density function $\rho(t,\cdot)$ of $(u,v)_t$ and the probability density function $p(t,\cdot)$ of $(y,z,v)$ satisfy
\[\rho(t,(u,v)) = p(t,(y,z,v)), \quad u = (y \odot y + z \odot z)^{1/2},\]
if the initial density function is cylindrical, i.e.
\[\rho(0,(u,v)) = p(0,(y,z,v)), \quad \text{ for all } u = (y \odot y + z \odot z)^{1/2}.\]

\section{Numerical Method for Cartesian Parametrisation}
We aim to numerically simulate the system (\ref{yzvproc}) and compare its performance to that of Hadamard-Langevin applied to (\ref{uvproc}). 
\subsection{Euler-Maruyama}
We first try the Euler-Maruyama scheme. Applying to (\ref{cyl}), we obtain the update, with stepsize $h$,
\[X_{n+1} = X_n - h\lambda X_n \, - h\nabla G_{3}(X_n)\, + \sqrt{2h/\beta}\,R_n,\]
for $R_n \sim N(0,I_{3d})$. However, due to the singularity $\frac{1}{u_t}$, this method may be unstable.

\subsection{Splitting Methods}
Suppose we have a system of differential equations as
\begin{equation}  
 \begin{aligned}
  \frac{d\vec x}{dt} = f(\vec x), \quad  x(0) = x_0,
  \end{aligned}\label{ivp_ode}
\end{equation}
and that $f$ is separable as 
\begin{equation}  
 \begin{aligned}
  f = f_1 + f_2 + \hdots + f_k, \quad k \geq 2,
  \end{aligned}\label{sep_f}
\end{equation}
so that each initial value problem 
\begin{equation}  
 \begin{aligned}
  \frac{d\vec x}{dt} = f_i(\vec x), \quad  x(0) = x_0,
  \end{aligned}\label{ivp_ode_small}
\end{equation}
is easier to solve than (\ref{ivp_ode}). The idea of splitting methods is that we take advantage of the decomposition (\ref{sep_f}) to approximate the solution to (\ref{ivp_ode}). ADD BCH STUFF, EXPLAIN WHY BENS METHOD WORKS.

%ADD INTRO/DESCRIPTION OF SPLITING METHODS HERE OR IN RELEVANT SECTION IN THESIS. DISCUSS

Our system can be tackled by splitting methods via dividing it into multiple parts and treating them sequentially. For example, we can divide the system (\ref{cyl}), in a similar fashion to that of the Langevin Dynamics splitting by Leimkuhler and Matthews \cite{leimkuhler_baoab}:
\begin{equation}  
 \begin{aligned}
 d\vec x= \underbrace{\Bigl[-\lambda \vec x \,dt +\sqrt{2/\beta}\, dW^3(t) \Bigr]}_{O} - \underbrace{\Bigl[ \nabla G_{3}(\vec x)\Bigr]}_{B},
  \end{aligned}\label{cyl_split}
\end{equation}
Part $O$ is an Orstein-Uhlenbeck system and can be solved exactly in one timestep. In a similar vein to the methods proposed by Leimkuhler et al. we propose the BOB scheme. Akin to Strang splitting scheme, we have
\begin{equation} 
\begin{aligned} 
(B) \quad &\vec x_{n+1/2}&=&\quad \vec x_{n} - \frac{h}{2}\,\nabla G_{3}(\vec x_n),\\
(O) \quad &\hat{\vec x}_{n+1/2} &=&\quad \vec x_{n+1/2}\, e^{-\lambda h} + \sqrt{\frac{1}{\beta\lambda} \left(1-e^{-2\lambda  h} \right)}\, \xi, \quad \xi \sim N(0,I_{3d}), \\
(B) \quad &\vec x_{n+1}&=&\quad \hat{\vec x}_{n+1/2} - \frac{h}{2}\,\nabla G_{3}(\hat{\vec x}_{n+1/2}),
\label{bob_split}
\end{aligned}
\end{equation}
Also try the sequential (BO) scheme, akin to Lie-Trotter for ODEs, should perform worse than BOB due to asymmetry creating $\mathcal{O}(h)$ terms in BCH expansion:
\begin{equation} 
\begin{aligned} 
(B) \quad &\vec x_{n+1/2}&=& \quad \vec x_{n} - h\,\nabla G_{3}(\vec x_n)\\
(O) \quad &\vec x_{n+1} &=& \quad \vec x_{n+1/2}\, e^{-\lambda h} + \sqrt{\frac{1}{\beta\lambda} \left(1-e^{-2\lambda  h} \right)}\, \xi, \quad \xi \sim N(0,I_{3d}),
\label{seq_split}
\end{aligned}
\end{equation}

We compare performance to the Euler-Maruyama scheme:
\begin{equation} 
\begin{split} 
\vec x_{n+1}&=\vec x_{n} - h \,\lambda \vec x_{n} - h\,\nabla G_{3}(\vec x_n) +\sqrt{2/\beta} \, R_n\, \quad R_n \sim N(0,I_{3d})
\label{euler}
\end{split}
\end{equation}
Additionally, by taking the limit of the friction coefficient $\gamma \to \infty$ in the BAOAB scheme for underdamped langevin dynamics, we can derive a new pseudo-Euler-Maruyama scheme, that has $\mathcal{O}(h^2)$ convergence for overdamped langevin dynamics \cite{leimkuhler_baoab}:
As found in BAOAB paper:
\begin{equation} 
\begin{split} 
\vec x_{n+1}&=\vec x_{n} - h \,\lambda \vec x_{n} - h\,\nabla G_{3}(\vec x_n) +\sqrt{2h/\beta} \, (R_n\, + R_{n+1}) \quad R_n, R_{n+1} \sim N(0,I_{3d})
\label{pseudo-euler}
\end{split}
\end{equation}
Above doesn't approximate invariant measure that well. From trial and error, found that
\begin{equation}
\begin{aligned}
\vec x_{n+1}&=\vec x_{n} - h \,\lambda \vec x_{n} - h\,\nabla G_{3}(\vec x_n) +\sqrt{h/2\beta} \, (R_n\, + R_{n+1}) \quad R_n, R_{n+1} \sim N(0,I_{3d})
\end{aligned}
\end{equation}
works pretty well and demonstrates $\mathcal{O}(h^2)$ convergence in test cases, as predicted by Leimkuhler.


\section{Numerics - 1D Toy Problem}
We simulate a toy problem. We use $G(x) = (Ax-y)^2$, with $A \in \mathbb{R}$, $y = Ax_0$, for $x_0 = -2$, $\lambda = 0.6 \times |Ay| = $. We simulate 10,000 particles, each for 10,000 iterations. We present our results below.






\begin{thebibliography}{99}
\bibitem{riskml}
Abadie, A. \& Kasy, M. Choosing among regularized estimators in empirical economics. {\em The Review of Economics and Statistics} (2019) 101 (5): 743–762.

\bibitem{leimkuhler_baoab}
 B. Leimkuhler \& C. Matthews. Rational construction of stochastic numerical methods for molecular sampling. {\em Applied Mathematics Research eXpress}, 2013(1):3456, 2013.
\end{thebibliography}

\end{document}
